\section{OOP - Classes and Objects}

\subsection{Section Overview}
Welcome to Object-Oriented Programming (OOP) in C++. OOP is a programming paradigm that uses ``objects'' to design applications and computer programs. It's a powerful way to think about and organize your code, especially for large, complex projects. In this section, we'll introduce the core concepts of OOP: classes and objects.

\subsection{What is Object-Oriented Programming?}
Object-Oriented Programming is based on several key principles:
\begin{itemize}
	\item \textbf{Encapsulation}: Bundling data (attributes) and methods (functions) that operate on the data into a single unit, or ``class''.
	\item \textbf{Abstraction}: Hiding complex implementation details and showing only the essential features of the object.
	\item \textbf{Inheritance}: Creating new classes from existing classes, inheriting their properties and behavior.
	\item \textbf{Polymorphism}: The ability of an object to take on many forms, most commonly through method overriding.
\end{itemize}
We will cover these concepts in detail throughout this part of the course.

\subsection{What are Classes and Objects?}
A \textbf{class} is a blueprint for creating objects. It defines a set of attributes (member variables) and methods (member functions) that the created objects will have. An \textbf{object} is an instance of a class. If a class is the blueprint for a house, an object is an actual house built from that blueprint.

\subsection{Declaring a Class and Creating Objects}
You declare a class using the \texttt{class} keyword. Inside the class, you can declare member variables and member functions. The \texttt{public} and \texttt{private} keywords control access to the members.
\begin{lstlisting}
	class Player {
	public:
		// attributes
		std::string name;
		int health;
		int xp;

		// methods
		void talk(std::string text) {
			std::cout << name << " says: " << text << std::endl;
		}
	};

	int main() {
		Player frank; // create a Player object (an instance of the Player class)
		frank.name = "Frank";
		frank.health = 100;
		frank.xp = 12;
		frank.talk("Hi there!");

		Player hero;
		hero.name = "Hero";
		hero.talk("Hello, Frank!");

		return 0;
	}
\end{lstlisting}

\subsection{Accessing Class Members}
Class members are accessed using the dot (\texttt{.}) operator for objects, and the arrow (\texttt{->}) operator for pointers to objects.
\begin{lstlisting}
	Player frank;
	frank.name = "Frank"; // using dot operator

	Player *enemy = new Player();
	enemy->name = "Enemy"; // using arrow operator
	enemy->talk("I will destroy you!");
	delete enemy;
\end{lstlisting}

\subsection{Implementing Member Methods}
Methods can be implemented inside the class declaration (inline) or outside. It is common practice to separate the declaration from the implementation.
\begin{lstlisting}
	class Account {
	private:
		double balance;
	public:
		void set_balance(double bal); // declaration
		double get_balance();       // declaration
	};

	// implementation using the scope resolution operator ::
	void Account::set_balance(double bal) {
		balance = bal;
	}
	double Account::get_balance() {
		return balance;
	}
\end{lstlisting}

\subsection{Coding Exercise 33: Add a method to an existing class}
Given the \texttt{Dog} class below, add a \texttt{bark()} method that prints ``Woof!''.
\begin{lstlisting}
	#include <iostream>
	#include <string>

	class Dog {
	public:
		std::string name;
		int age;
		void bark() {
			std::cout << "Woof!" << std::endl;
		}
	};

	int main() {
		Dog my_dog;
		my_dog.name = "Fido";
		my_dog.age = 5;
		my_dog.bark();
		return 0;
	}
\end{lstlisting}

\subsection{Separating Specification from Implementation}
For better organization, it is common to put the class declaration in a header file (\texttt{.h} or \texttt{.hpp}) and the implementation of the member functions in a source file (\texttt{.cpp}).

\subsection{The include Guard}
To prevent a header file from being included multiple times in the same source file (which would cause redefinition errors), we use include guards.
\begin{lstlisting}
	// In Account.h
	#ifndef ACCOUNT_H   // Note: avoid leading underscores — reserved by the standard
	#define ACCOUNT_H

	class Account {
		// ...
	};

	#endif // ACCOUNT_H

	// A more modern, widely supported (but not officially standardized) alternative:
	#pragma once
	class Account {
		// ...
	};
\end{lstlisting}

\subsection{Constructors and Destructors}
A \textbf{constructor} is a special member function that is automatically called when an object of a class is created. It is used to initialize the object's attributes. A \textbf{destructor} is automatically called when an object is destroyed. It is used to release resources.

\subsection{The Default Constructor}
A default constructor is a constructor that takes no arguments. If you don't provide any constructors, the compiler will generate a default constructor for you.
\begin{lstlisting}
	class Player {
	public:
		// Preferred: member initializer list
		Player() : name{"None"}, health{100}, xp{0} {}

		// Also valid but less efficient: assignment in body
		// Player() { name = "None"; health = 100; xp = 0; }
		// ...
	};
	Player frank; // Default constructor is called
\end{lstlisting}

\begin{remark}
	Prefer \textbf{member initializer lists} (\texttt{: name\{"None"\}, ...}) over assignment in the constructor body. Initializer lists construct members directly with the right value, while the body-assignment approach default-constructs them first and then assigns — which is redundant and slower for types like \texttt{std::string}.
\end{remark}

\subsection{Overloading Constructors}
You can have multiple constructors with different parameter lists. This is constructor overloading.
\begin{lstlisting}
	class Player {
	public:
		Player(); // Default constructor
		Player(std::string name_val); // Overloaded constructor
		Player(std::string name_val, int health_val, int xp_val); // Overloaded
	};

	Player frank; // Player()
	Player hero("Hero"); // Player(std::string)
	Player villain("Villain", 100, 55); // Player(std::string, int, int)
\end{lstlisting}

\subsection{Choosing a Constructor}
The compiler chooses which constructor to call based on the arguments you provide when creating the object.

\subsection{Default Constructor Parameters}
You can use default parameter values in your constructors to create a more flexible ``default'' constructor.
\begin{lstlisting}
	class Player {
	public:
		Player(std::string name_val = "None", int health_val = 100, int xp_val = 0);
	};

	Player frank; // Player("None", 100, 0)
	Player hero("Hero"); // Player("Hero", 100, 0)
\end{lstlisting}

\subsection{Coding Exercise 34: Adding a Default Constructor to an existing class}
Add a default constructor to the \texttt{Car} class that initializes \texttt{make} to ``Unknown'' and \texttt{year} to 0.
\begin{lstlisting}
	#include <iostream>
	#include <string>

	class Car {
	public:
		std::string make;
		int year;
		Car() {
			make = "Unknown";
			year = 0;
		}
	};

	int main() {
		Car my_car;
		std::cout << "Make: " << my_car.make << ", Year: " << my_car.year << std::endl;
		return 0;
	}
\end{lstlisting}

\subsection{The Copy Constructor}
A copy constructor is used to create a new object as a copy of an existing object.
\begin{lstlisting}
	class Player {
	public:
		Player(const Player &source); // Copy constructor prototype
	};

	Player frank("Frank", 100, 10);
	Player copy_of_frank {frank}; // Copy constructor is called
\end{lstlisting}

\subsection{Shallow vs. Deep Copying}
When a class has pointer member variables, the default copy constructor performs a \textbf{shallow copy} — it copies the pointer's value, not the data it points to. This means two objects end up sharing the same heap memory, causing a double-free crash when both destructors run.

\begin{example}[Shallow Copy Problem and Deep Copy Fix]{\leavevmode}
\begin{lstlisting}
	class Shallow {
		int *data;
	public:
		Shallow(int d) : data{new int{d}} {}
		~Shallow() { delete data; }
		// Compiler-generated copy constructor does:
		//   data = other.data;  <-- both point to the same int!
	};

	Shallow a{42};
	Shallow b{a}; // b.data == a.data  => double-free when both destruct!

	class Deep {
		int *data;
	public:
		Deep(int d) : data{new int{d}} {}
		~Deep() { delete data; }
		// User-defined copy constructor allocates its own storage
		Deep(const Deep &other) : data{new int{*other.data}} {}
	};

	Deep c{42};
	Deep d{c}; // d.data is a separate copy — safe
\end{lstlisting}
\end{example}

\begin{remark}
	\textbf{Rule of Three:} If your class needs to define any one of (1) a destructor, (2) a copy constructor, or (3) a copy assignment operator, it almost certainly needs all three. This rule exists because all three are needed to safely manage a raw pointer resource. In C++11 and later this extends to the \textbf{Rule of Five}, adding the move constructor and move assignment operator.
\end{remark}

\subsection{The move Constructor}
C++11 introduced the move constructor, which is used for transferring ownership of resources from one object to another. It's an optimization that can avoid expensive deep copies.

\subsection{The this Pointer}
Within a member function, the \texttt{this} pointer is a pointer to the object that the method was called on. It allows you to access the object's own members.
\begin{lstlisting}
	void Player::set_name(std::string name) {
		this->name = name; // 'this->name' is the member, 'name' is the parameter
	}
\end{lstlisting}

\subsection{Using const with Classes}
You can have \texttt{const} member functions that are not allowed to modify the object's state. This is good practice for ``getter'' methods.
\begin{lstlisting}
	class Player {
	public:
		std::string get_name() const; // A const member function
	private:
		std::string name;
	};

	std::string Player::get_name() const {
		// health = 100; // This would cause a compiler error
		return name;
	}
\end{lstlisting}

\subsection{static Class Members}
A \texttt{static} member variable is shared by all objects of the class. A \texttt{static} member function can be called without creating an object of the class, and it can only access \texttt{static} member variables.
\begin{lstlisting}
	class Player {
	public:
		static int num_players;
		Player() { num_players++; }
		~Player() { num_players--; }
	};
	int Player::num_players = 0; // Initialize static variable

	int main() {
		Player frank;
		Player hero;
		std::cout << "Num players: " << Player::num_players << std::endl; // 2
	}
\end{lstlisting}

\subsection{struct vs class}
In C++, a \texttt{struct} is almost identical to a \texttt{class}. The only difference is that members of a \texttt{struct} are \texttt{public} by default, while members of a \texttt{class} are \texttt{private} by default. By convention, \texttt{struct}s are used for simple data aggregates without complex methods.

\subsection{Friends of a Class}
A \texttt{friend} function or \texttt{friend} class can access the \texttt{private} and \texttt{protected} members of a class. This should be used sparingly, as it can break encapsulation.

\subsection{Section Challenge}
Create a \texttt{Movie} class with attributes for name, rating (e.g., ``PG-13''), and a watched count. Implement a constructor, a copy constructor, a destructor, and methods to increment the watched count and display the movie's details. In \texttt{main}, create a vector of \texttt{Movie} objects and test your methods.
\subsection{Section Challenge --- Solution}
\begin{lstlisting}
	#include <iostream>
	#include <string>
	#include <vector>

	class Movie {
	private:
		std::string name;
		std::string rating;
		int watched_count;

	public:
		Movie(std::string name, std::string rating, int watched)
			: name(name), rating(rating), watched_count(watched) {
			std::cout << "Constructor called for " << name << std::endl;
		}

		Movie(const Movie &source)
			: Movie(source.name, source.rating, source.watched_count) {
			std::cout << "Copy constructor called for " << name << std::endl;
		}

		~Movie() {
			std::cout << "Destructor called for " << name << std::endl;
		}

		void increment_watched() { watched_count++; }

		void display() const {
			std::cout << name << ", " << rating << ", " << watched_count << std::endl;
		}
	};

	void add_movie(std::vector<Movie> &movies, const std::string &name, const std::string &rating, int watched) {
		movies.push_back(Movie(name, rating, watched));
	}

	int main() {
		std::vector<Movie> my_movies;
		add_movie(my_movies, "The Matrix", "R", 5);
		add_movie(my_movies, "Inception", "PG-13", 3);

		my_movies.at(0).increment_watched();
		my_movies.at(0).display();

		return 0;
	}
\end{lstlisting}

\subsection{Quiz 10: Section 13 Quiz}
\begin{enumerate}
	\item What is the main purpose of a constructor?
	\begin{enumerate}
		\item[a)] To destroy an object.
		\item[b)] To initialize an object's member variables.
		\item[c)] To copy an object.
		\item[d)] To define the class.
	\end{enumerate}

	\item What is the difference between a \texttt{class} and a \texttt{struct} in C++?
	\begin{enumerate}
		\item[a)] They are completely different.
		\item[b)] \texttt{class} members are \texttt{public} by default; \texttt{struct} members are \texttt{private} by default.
		\item[c)] \texttt{class} members are \texttt{private} by default; \texttt{struct} members are \texttt{public} by default.
		\item[d)] There is no difference.
	\end{enumerate}

	\item A member function declared with \texttt{const} at the end:
	\begin{enumerate}
		\item[a)] Can modify the object's member variables.
		\item[b)] Cannot modify the object's member variables.
		\item[c)] Is a static function.
		\item[d)] Must be public.
	\end{enumerate}
\end{enumerate}

\textbf{Answers:} 1. (b), 2. (c), 3. (b)
